\chapter{Photoeffects}

The literature on photoeffects is dominated by two publications. First was a
book by Kinzel which reported studies on 970 species (ref. 17). These were grouped
into those in which germination was favored by light, those in which germination was
not favored by light, and those in which light had no effect. The second publication was
Borkwith and coworkers (ref. 18) which demonstrated that the effectsof light were
reversible in lettuce with implications that photoeffects in germination of seeds were
related to the phytochrome equilibrium. A brief review of these two publications is
given at the end of this chapter.

The responses of germination to light were \textit{complex}. Usually light was either
required for germination, had no effect, or blocked germination. This differs from the
past literature which described the germinations as "favored or unfavored" by light. It
is suspected that in the past the presence or absence of light was not rigidly controlled.
An entirely unexpected result was that many of the species that required light for
germination would also germinate in the dark if treated with gibberelic acid-3. This
subject is deferred until the next chapter.

The requirement of light for germination was characteristic of swamp plants. It
also occurred frequently in woodland plants and in many Ericaceae and Solanaceae.
In swamps and woodlands light is more apt to be a limiting factor than moisture so the
requirement of light for germination has survival value. The response to light was
sometimes affected by prior treatments such as dry storage or three months moist at
40, but in some species pretreatments had no effect on the response to light. When
seed catalogs say surface sow, it can be inferred that light is required.

The genus Asclepias demonstrates a range of behaviors in a single genus.
Light has no effect on the germination of dry land species such as Asclepias tuberosa
whereas it is an absolute requirement for germination in a swamp species such as
Asclepias incamata. Species from an intermediate evironment such as Asclepias
syriacus have an intermediate behavior. The example of Asclepias illustrates the
general trend that photorequiremerits are more associated with the environment in
which the species grows than with the family or genera. Thus the photorequirement
was found in both monocotyledons and dicotyledons and in both epigeal and
hypogeal germinators. This suggests that the mechanism of the photorequirement
must be very old in an evolutionary sense and that all families have some common
chemical system that can be activated to generate a photorequirement for germination.
The responses to light can be divided into seven categories, and each are now
discussed.

It was convenient to divide the behaviors into seven categories. \textit{The first
category} consists of species which germinated only in light at 70 and the germination
rates and percent germinations were relatively insensitive to prior treatments. Thus six
months of dry storage or up to a year of alternating moist cycles at 40 and 70 had little
or no effect. This group uses the photorequirement as the 2ft mechanism for
preventing germination of the seeds before they are dispersed. A complete set of
experiments were conducted in most, but not all, of the following so that future
experiments may lead to some revision. Some members of this category are listed in

Table 11-1. Species That Use Only a Photorequirement to Delay Germination

\begin{multicols}{2}
Aconitum heterophyllum

Aletris farinosa

Angelica grayi

Asciepias phytolaccoides

Astilbe chinensis

Campanula takesimana

Campsis radicans

Cephalanthus occidentalis

Chelone glabra

Clematis verticillata

Clematis virginiana

Dlplacus bifidus

Dipsacus sylvestris

Epigea repens

Gaultheria trichophylla

Gentiana saponaria

Geum coccineum

Helonias bullata

Hydrangea quercifolia

Iris milesii

Iris pseudoacorus

Iris sibirica

Iris spuria

Iris tectorum

Iris tectorum alba

Kalmia angustifolia

Kalmia latifolia

Lamium maculatum

Lobelia cardinalis

Lobelia syphilitica

Lyonia mariaria

Macleaya cordata

Meconopsis aculeata

Mimulus primuloides

Mimulus ringens

Monarda fistulosa

Nepeta cataria

Oenothera biennis

Pinguicula macroceras

Peltoboykinia tellimoides

Penstemon digitalis

Penstemon frutescens

Penstemon hirsutus

Penstemon hirsutus pygmaea

Penstemon smallil

Physelis alkekengi

Physalis virginiana

Polemonium grayanum

Polemonium vanbruntiae

Ranunculus repens

Rhexia mariana

Rhexia virginica

Rhododendron lepidotum

Spirea betulaefolia

Spirea japonica

Telesonix jamesii

Tiarella cordifolia

Tiarella wherryi

Tricyrtis affinis

Tricyrtis dilitata

Verbascum blattaria

Viola tricolor
\end{multicols}

The second category consists of species which not only require light to
germinate but also a prior pretreatment. Leucothoe racemosa and Primula kisoana
required six months dry storage at 70 before they would germinate in light. Asclepias
incarnata, Campanula americana, and Verbena hastata required three months moist
at 40 before they would germinate in light.

The third category consists of species whose germination was totally blocked
by light. A prior three months at 40 or a prior three months in light at 70 had little or no
effect. On shifting to the dark at 70, germination began. The induction time, the
germination rate, and the percent germination were the same as if the prior treatment
had not occurred. Examples of such species are Convallaria majalis, Cyclamen
persicum, Ophiopogon japonicus, and Poncirus trifoliata: Schizanthus approximated
this behavior. The earlier experiments were notdesigned to recognize this behavior,
and there may be more examples of this behavior than the meagre list would suggest.

The fourth category consists of species that absolutely require light for
germination of fresh seed at 70, but various pretreatments remove this
photorequirement. In some such as Paulownia tomentosum and Silene armeria the
photorequirement is eliminated by six months dry storage at 70. The Paulownia was
particularly interesting because fresh seed germinated nearly 100\% in 70L and gave
not a single germination in 70D in samples of 500 seeds. After two years of dry
storage at 70, the seed germinated in either 70L or 70D. The induction times and -
germination rates and the percent germination (100\%) were identical in either the light
(70L) or the dark (70D). It is possible that more species were in this category, but were
not recognized because only DS seed was available for study.

The behaviors of Ailanthus altissima, Gentiana scabra, and Solanum
dulcamara were more complex, but the essence was that the photo requirement was
eliminated by either three months at 40 or by dry storage. The effects of the chilling
period and of the dry storage were additive suggesting that both are destroying the
same chemical inhibitor(s) and that these are the same inhibitor(s) destroyed by light.
The fifth grou p consists of species that required light for germination at 70, but
they also germinated in the dark at 40. The percent germination at 40 is lower and the
germination rates are slower, but percent germination is enough to constitute
satisfactory germination. Some examples are Calyptridium umbellatum, Centaurium
scilloides, Cynoglossum amabile, Gentiana autumnalis, Geum montanum, Tricyrtis
hirta, Tricyrtis hirta alba, and Tricyrtis macrantha.

The sixth category consists of species that germinate in either light at 70 (70L)
or in outdoor conditions. Germination fails in all other treatments.

The seventh category consists of species whose germination is promoted by
light, but some germination occurs in the dark. Members of this group are listed in
Table 11-2. This category fits the favored by light description of Kinzel (ref. 17). It is
possible that these actually have an absolute photorequirement, but only a very short
exposure to light is required. Thus the short exposure to light that occurs in making the
observations might have been enough to trigger some germination. It is also possible
that some of this group showed a slightly higher germination in light at 70 because
absorption of light caused the temperature of the experiment in light to be slighly
above the temperature of the dark experiments. This possibility was tested with
Gomphrena and found to be the explanation (see Gomphrena in Chapter 20). This
possibility was also examined by placing seeds of Lobelia cardinalis on different
colored paper and placing under the lights. The idea was that the darker the paper the
higher the effective temperature with the black paper being the warmest and the white
paper the closest to the room temperature. The seeds gave the same germination on
all papers.

Table 11-2. Species Whose Germination Is Increased by Light but Light Is Not Essential

\begin{multicols}{2}
Campanula punctata

Campanula ramosissima

Campanula rapunculoides

Catalpa bignonioides

Cichorum intybus

Clematis connata

Clematis grata

Datura stramonium

Hypericum perforatum

Itea virginica

Lisianthus russelianum

Lychnis wilfordii

Lythrurn salicaria

Parnassia fimbriata

Parnassia glauca

Pieris japonica

Platanus occidentalis

Primula pamirica

Primula warshenewskiana

Ranunculus adoneus

Rhododendron cawtawbiense

Sedum puichellum

Spraguea umbellata

Symplocarpus foetidus

Trichostema setaceum
\end{multicols}

Two curious effects were observed with the aquatic Hymenocallis occidentalis
and Zantedeschia aethiopica. The Hymenocallis has a thick green seed coat. It is
necessary to keep this on and to keep it photosynthesizing in light in order to get
germination several months later. Presumable the photosynthesis of the seed coat
produces chemicals essential for the ultimate germination. In the Zantedeschia there
is a green photosynthesizing seed capsule, and this must be removed and the seeds
washed to give germination. Presumably the seed capsule is producing an inhibitor
by its photosynthesis.

This brings us to commenting on the past literature. The Kinzel work (ref. 17) is
always summarized (refs. 2-7) in terms of germination being favored or not favored.
Perhaps Kinzel is being misquoted or perhaps something is lost in translation or
perhaps Kinzel conducted the experients in media where the variable of tight was not
well controlled. In any event the word favored suggests minor effects. Unfortunately, I
have not been able as yet to get a copy of Kinzels book (it was published in german in
1926), and the summaries of his work are so qualitative as to make one wonder if the
reviewers had read the original work. At least no critical analysis has appeared. In the
present work, although the term favored is possibly applicable to the species in Table
11-2, it was more common for light to be an absolute requirement for germination or to
absolutely block germination.

The Borkwith work (ref. 18) is cited (refs. 1-6) as the classic demonstration of
photoeffects and the demonstration that such effects can be reversible. Specifically
fresh seeds of Grand Rapids lettuce were irradiated with alternating cycles of one
minute of red light and three minutes of infrared light. Germination was 10\% for the
dark blank, 90\% if the seed had been last exposed to red light, and 50\% if the seeds
had been last exposed to infrared light. The Borkwith paper led to many studies
designed to determine whether or not the phytochrome photoequilibria, so important in
photosynthesis, was responsible for the effects in seed germination. A recent book on
phytochrome (ref. 19) has a chapter by J. D. Smith on seed germination, and the
author concludes that the question is still unsettled. The Borkwith papers also inspired
studies on varying light periods. It has been reported that different strains of Zea mays
have different responsesto light (ref. 20) and that differing light cycles had different
effects on germination of Kalanchoe (ref. 21). None of this seems of much practical
importance for the successful germination of seeds.

The experiments and conclusions of Borkwith and his coworkers disturb me.
First, measurement of induction periods and rates would have shown more than the
percent germinations that were reported. Second, my experiments with Paulownia
show that several days of exposure to light were required to reach the maximum
germination (see Paulownia in Chapter 20). Third, in my view light is acting by
destruction of an inhibitor system. If so, the effects will be largely governed by the
wave lengths of the light, the intensity of the light, and the duration of exposure. Four,
the difference of 50\% and 90\% is minor compared to the absolute light requirements
and absolute blocks of germination by light which were characteristic of my own work.
Finally I have heard that the light effects on Grand Rapids lettuce rapidly disappears
on dry storage.

